\subsection{SISTEMAS DISTRIBUÍDOS E A INTERNET}

\begin{comment}
Adicionar:
  - Relacionar contexto do ISAM com o contexto da INTERNET

Melhorar:
  - Pegar funcionalidades do CORBA quando listar as funcionalidades de middlewares
  - Referência para embasar a afirmação sobre os componentes dos sitemas distribuídos

Estrutura do paragrafo
  - Contextualizar a internet como um conjunto de sistemas distribuídos
  - Papel dos protocolos na internet
  - Usar exemplos de aplicações reais, como editores compartilhados
  - Como peer-to-peer se relaciona com a internet (DNS, IPv4, IPv6)
  - Indicar as relações dos servidores root com a comunicação peer-to-peer
\end{comment}


A \en{internet} ocupa um papel essencial na atualidade, através da computação móvel, aplicações empresariais, redes \en{LAN}, redes sem-fio (\en{WIFI}), e outros exemplos de conectividade que possibilitam a conexão. Essas conexões ocorrem de diversas formas: \en{WiMAX} (interoperabilidade mundial de acesso a micro-ondas), \en{Bluetooth}, quarta geração de redes móveis (4G) e outros exemplos de trammissão de dados pelo meio físico. Para possibilitar essa comunicação há os protocolos de rede como o \en{TCP} e o \en{IP} que, dentre outros, formam o modelo referêncial \en{TCP/IP} que define a pilha de protocolos reponsáveis por ligar os dados vindos pelos meios físicos às aplicações \cite{alberti2016distributedsystemsfutureinternet}. Nota-se, desse modo, como sistemas distintos conseguem cooperar na \en{internet}.

Esse conjunto de sistemas distintos dividem a tarefa de processamento de dados entre si, formando sistemas distribuídos. Um sistema distribuído é um ambiente de computação, onde vários componentes localizados em multiplos dispositivos de computação em uma rede segregam trabalho entre si, de modo que, para o usuário final, pareça ser um sistema único que realiza o trabalho \cite{khole2023compendium}. Os componentes de um sistema distribuídos são \en{softwares} com responsabilidades específicas, como: \en{APIs} que são resposáveis por conectar o cliente com o componente de persistência de dados e aplicar a lógica de negócio; \en{middlewares} responsáveis por afunilar o trafégo de dados, filtrando dados, aplicando restrição de acesso, validações em requisições, etc. 
